\section{优化方向与异步化路径选择}\label{sec:opt-motivation}\label{sec:async-feasibility}

本章优化的直接目标是提高在线推理过程的有效 FPS。第三章已经表明，
在树莓派 5 上，单纯依靠同步流程难以达到 30 FPS；因此需要从控制循环结构入手，
寻找可以从主路径中移出的阻塞环节。相机读取已经采用后台线程异步获取最新帧，
主循环只取最近一次结果。这一机制说明，在不改变任务语义的前提下，
部分 I/O 或计算过程可以提前执行，从而减少主循环等待时间。

基于这一思路，本文首先考察观测、推理和执行三个阶段中可异步化的事件。
舵机读写看似也可以交给后台任务，但 SO-101 机械臂的多个舵机共享同一条
TTL 半双工总线，同一时刻只能进行一次读或写。为了避免后台写入和主线程读状态
同时占用总线，必须在总线访问外层加互斥锁；一旦加锁，读写仍然只能按顺序执行，
异步线程只是把原来的串行访问换成“持锁串行 + 线程切换”，
并不能真正隐藏总线时间。舵机写入还涉及动作限幅的连续性，延迟到后台执行后，
错误反馈和安全判断都会变得滞后。

舵机读取同样不能直接照搬相机的后台线程模式。若增加一个后台读线程，
它仍然要访问同一条舵机总线，并会与主线程写目标位置竞争同一把总线锁；
锁住之后，读写事务仍按顺序通过总线，异步读无法形成真正并行。
一种较温和的方案是采用流水线读：本帧接收上一帧请求的状态，
同时发出下一帧请求。该方案可以避免后台线程和总线锁竞争，
但代价是观测天然滞后一帧，并且仍需改动底层舵机通信封装。
第三章测得单次舵机读取约 1.2 ms，相对 ACT 推理约 2 s 的主瓶颈而言收益过小，
因此不作为本章主要实验方向。

相比之下，推理阶段既是主要耗时来源，又天然具备与动作执行重叠的时间窗口。
在 chunk\_size 为 100、目标 30 FPS 的配置下，一个 action chunk
执行约 3.3 s；同步模式在这段时间内并不会提前计算下一段动作。
因此，将下一次模型推理移出主控制路径，使其在当前 chunk 执行期间并行进行，
是更符合当前瓶颈结构的异步化方向。后续 4.2 节分析 LeRobot 的双进程推理机制，
4.3 节进一步通过异步推理触发阈值扫描验证该方案的实际收益和边界。
